\documentclass[11pt, a4paper]{article}

% --- Packages ---
\usepackage[utf8]{inputenc}
\usepackage[T1]{fontenc}
\usepackage{geometry}
\geometry{margin=1in}
\usepackage{hyperref}
\usepackage{graphicx}
\usepackage{booktabs}
\usepackage{authblk}
\usepackage{fancyhdr}

% --- Page Style ---
\pagestyle{fancy}
\fancyhf{}
\renewcommand{\headrulewidth}{0pt}
\fancyfoot[C]{Page \thepage\ of 13}

% --- Title Information ---
\title{\textbf{fmsg} \\ \large A message definition and protocol}
\author{Mark Mennell \\ \href{mailto:markmnl@fmsg.io}{markmnl@fmsg.io}}
\date{2026}

\begin{document}

\maketitle

\begin{abstract}
A message definition and protocol for sending and receiving messages. Compared to email's SMTP and IM applications, fmsg has many overlapping features; like email, fmsg is an open protocol anyone can host, unlike email but like IM applications fmsg is efficient at conversational message threads. Included in the fmsg design are security features built-in such as sender verification and message integrity. Sent via a fmsg host to one or more recipients, each message in a thread is linked to the previous using a cryptographic hash. Overall fmsg provides an efficient and secure messaging system with ownership and control at the domain level.
\end{abstract}

\newpage
\tableofcontents
\newpage

\section{Introduction}
Electronic mail (email) has co-existed with The Internet since inception and has become ubiquitous for messaging, from personal to corporate and even official communications. Meanwhile Instant Messaging (IM) applications have proliferated with businesses building products with ever increasing features. Many such applications have become part of everyday life. Despite the advancement of IM applications, email still maintains a stronghold for certain types of communication.

This paper explores successes and issues with the everyday messaging applications we use: email and IM applications. fmsg (thought of as "f-message", always stylised lower-case) is then introduced a binary protocol intended as an alternative for both. Feature wise comparison is made between SMTP, IM applications and fmsg. Real world comparison is made with an fmsg implementation: fmsgd. Appended is the fmsg specification and links to relevant code and document repositories.

\section{Stronghold of Email}
Email remains widely used today because it offers ownership and control at the domain level: messages travel via Domain Name Service (DNS) defined servers, server-to-server communication, allowing domain owners especially corporations to manage their own servers and meet requirements around privacy, data sovereignty, cost, and legal obligations. Since nobody owns email itself, anyone can run their own email service, while personal users overwhelmingly rely on hosted providers to avoid the complexity of managing infrastructure.

Email also carries a sense of formality due to its longer, richer formatting options and its delayed, non-conversational nature. Perhaps another reason for email's persistence is an email address often forms the basis of an online identity online services use email for verification and communication with a user creating an account. So to use many online services you must have an email address first!

\section{Pain of Email}
The success and massive adoption of email over decades has inevitably led to issues, which fmsg aims to solve, namely; spam, efficiency, operating complexity and user experience - particularly with back and forth conversation. As of writing spam, or "junk" email, constitutes nearly half of all emails sent (debounce.com Email Spam Statistics 2025), which apart from the annoyance, consumes computing resources especially in filtering out junk.

Email uses the Simple Mail Transfer Protocol (SMTP) for server-to-server transmission, with clients also using SMTP to send messages, then protocols like Post Office Protocol (POP) and Internet Message Access Protocol (IMAP) to retrieve them; however, since SMTP sends each email standalone, replies typically include the entire original email content, causing long email "chains," or "threads," to grow quickly in size and making email unsuitable for conversational style communication. 

User Experience suffers because the practice of including original content in replies lacks a formalised standard. Attachments in email are sent as text! This alone increases the size of binary attachments by one third of their original size because they are encoded as text (base 64) and appended to the message.

\section{Instant Messaging Applications}
Along with the rise in online connected devices, IM apps have proliferated as a way to electronically message someone else. IM apps tend to send their messages immediately to central server(s) owned and controlled by the companies that made them.

\begin{table}[h]
\centering
\caption{Popular Instant Messaging Applications}
\begin{tabular}{@{}lll@{}}
\toprule
\textbf{Instant Messaging App} & \textbf{Approx. Monthly Active Users*} & \textbf{Company/ Owner} \\ \midrule
WhatsApp & 0 & 0 \\
WeChat (Weixin) & 0 & 0 \\
Facebook Messenger & 0 & 0 \\
Telegram & 0 & 0 \\
Snapchat & 0 & 0 \\
QQ & 0 & 0 \\ \bottomrule
\end{tabular}
\end{table}

\section{Features of fmsg}
Perhaps the best way to introduce fmsg is a feature level comparison to email and IM apps.

\subsection{Attached files}
Arbitrary files can be included with messages. Email and fmsg include arbitrary file attachments, some IM apps do too though often limited types.

\subsection{Audio and video calls}
Live audio and video transmission, only possible in IM apps usually using another protocol than the messaging one but initiated and bundled within the application.

\subsection{Message integrity*}
Are messages tamper proof in transit and at rest? Email itself has no integrity guarantees but can be configured using additional mechanisms such as TLS, DKIM, and DMARC. fmsg has message integrity built-in such that even after delivery, messages can be verified they are unaltered.

\subsection{Ownership and control}
Both email and fmsg are owned and controlled at the domain level, whereas IM applications are owned and controlled by their owner.

\subsection{Prone to spam}
Email is prone to spam. Any service that allows unsolicited messages inherently is prone. Email is particularly prone because such use was not considered in the base protocol (SMTP) design.

\subsection{Rich text formatting}
Practically email and fmsg support a subset of HTML for rich text formatting. IM applications support plain text with varying levels of formatting.

\section{Tour of fmsg}
A message consists of header fields (size, topic, type etc.) including a cryptographic digest (SHA-256) of a parent message, if any. Messages are sent via a sender's fmsg host to one or more recipients. As message threads, or "chains", are created, a hierarchical structure builds up with each reply a direct acyclic graph (DAG).

\section{The Challenge}
A defining feature of the fmsg protocol is the automatic challenge. The challenge is a separate connection opened by the receiving host to the sender during receipt of a message before the incoming message content will be downloaded and accepted.

\section{Addresses}
An fmsg address looks like an email address except has a leading "@" character: \texttt{@username@example.com}.

\section{Practical Considerations}
fmsg itself only describes host-to-host communication, leaving other functionality such as retrieval of messages, user identity and access management, undefined.

\section{Operational Comparison}
We compare a fully functional setup of fmsg (using fmsgd) vs. email (using mailcow). Size of message transmitted and time taken to send are recorded over a variety of scenarios.

\section{Method}
All email and fmsg tests are conducted over two setups: 1) sender and receiver(s) are on the same machine, and 2) sender and receiver(s) are on geographically distant machines.

\section{Results}
TODO

\section{Conclusion}
TODO

\section{Notes on Encryption}
The fmsg protocol itself does not encrypt messages – instead, much like email, leaves it up to user/client to encrypt message content. The whole message should be encrypted in transit using PKI such as TLS.

\appendix
\section{Appendix I - Links}
Documentation repository on GitHub: \url{https://github.com/markmnl/fmsg}

\section{Appendix II - Specification}
[Specification details would follow here]

\end{document}